% --- CORRECCIÓN IMPORTANTE EN LA PRIMERA LÍNEA ---
% Agregamos "xcolor=table" para que funcione \rowcolor sin errores
\documentclass[aspectratio=169, 10pt, xcolor=table]{beamer}

% --- Configuración del Tema (estilo profesional como el ejemplo) ---
\usetheme{Madrid}
\usecolortheme{dolphin}

% --- Paquetes necesarios ---
\usepackage[utf8]{inputenc}
\usepackage[spanish]{babel}
\usepackage{amsmath, amssymb}
\usepackage{booktabs}
\usepackage{graphicx}
\usepackage{listings}
\usepackage{tikz}
\usetikzlibrary{arrows.meta, positioning, shapes.geometric}

% --- Ruta de Imágenes (crea la carpeta "imagenes" y coloca ahí tus archivos) ---
\graphicspath{{./imagenes/}}

% --- Estilo para código Python ---
\definecolor{codegreen}{rgb}{0,0.6,0}
\definecolor{codegray}{rgb}{0.5,0.5,0.5}
\definecolor{codepurple}{rgb}{0.58,0,0.82}
\definecolor{backcolour}{rgb}{0.95,0.95,0.92}
\lstset{
    language=Python,
    backgroundcolor=\color{backcolour},
    commentstyle=\color{codegreen},
    keywordstyle=\color{blue},
    numberstyle=\tiny\color{codegray},
    stringstyle=\color{codepurple},
    basicstyle=\ttfamily\scriptsize,
    breaklines=true,
    frame=single,
    showstringspaces=false
}

% --- Pie de página personalizado ---
\makeatletter

\setbeamertemplate{footline}{
  \leavevmode%
  \hbox{%
  \begin{beamercolorbox}[wd=.333333\paperwidth,ht=2.25ex,dp=1ex,center]{author in head/foot}%
    \usebeamerfont{author in head/foot}\insertshortauthor
  \end{beamercolorbox}%
  \begin{beamercolorbox}[wd=.333333\paperwidth,ht=2.25ex,dp=1ex,center]{title in head/foot}%
    \usebeamerfont{title in head/foot}Sesión 2
  \end{beamercolorbox}%
  \begin{beamercolorbox}[wd=.333333\paperwidth,ht=2.25ex,dp=1ex,right]{date in head/foot}%
    \usebeamerfont{date in head/foot}NLP \hfill \insertframenumber/\inserttotalframenumber \hspace*{0.3cm}
  \end{beamercolorbox}}%
  \vskip0pt%
}

\makeatother

% --- Datos de la portada ---
\title[SESIÓN 2 NLP]{\textbf{Sesión 2}}
\subtitle{Preprocesamiento de Texto: \\ De texto crudo a datos listos para modelar}
\author[Prof. C. Sánchez]{Carlos César Sánchez Coronel}
\institute{\textbf{NLP}}
\date{}

% Logo opcional (descomentar si tienes la imagen)
% \titlegraphic{\includegraphics[height=1.5cm]{logo.png}}

% --- Eliminar la barra de navegación (opcional) ---
\setbeamertemplate{navigation symbols}{}

% --- Índice al inicio de cada sección ---
\AtBeginSection[]{
  \begin{frame}
    \frametitle{Contenido}
    \tableofcontents[currentsection]
  \end{frame}
}

% ============================================================================
% INICIO DEL DOCUMENTO
% ============================================================================
\begin{document}

% --- Portada ---
\begin{frame}
    \titlepage
\end{frame}

% --- Índice general ---
\begin{frame}{Contenido}
    \tableofcontents
\end{frame}

% ============================================================================
% SECCIÓN 1: INTRODUCCIÓN AL PIPELINE DE PREPROCESAMIENTO
% ============================================================================
\section{Introducción al Pipeline de Preprocesamiento}

\begin{frame}{¿Por qué preprocesar el texto?}
    \begin{itemize}
        \item El texto real es \textbf{ruidoso}: URLs, menciones, emojis, mayúsculas, puntuación, etc.
        \item El preprocesamiento transforma el texto crudo en una representación estructurada para modelos.
        \item Un pipeline típico incluye:
        \begin{enumerate}
            \item Limpieza inicial
            \item Tokenización
            \item Normalización
            \item Eliminación de stopwords (opcional)
            \item Reducción morfológica (stemming/lematización)
        \end{enumerate}
    \end{itemize}
    % Basado en introducción del PDF
\end{frame}

% ============================================================================
% SECCIÓN 2: TOKENIZACIÓN
% ============================================================================
\section{Tokenización}

\begin{frame}{Tokenización por palabras}
    \begin{block}{Definición}
        Dividir un texto en unidades mínimas (tokens), normalmente palabras, signos de puntuación, etc.
    \end{block}
    \begin{itemize}
        \item \textbf{Desafíos}: contracciones (``don't''), puntuación, idiomas sin espacios.
        \item \textbf{Herramientas}: NLTK, spaCy, Stanford CoreNLP.
    \end{itemize}
    % Basado en PDF
\end{frame}

\begin{frame}{Tokenización por oraciones}
    \begin{block}{Definición}
        Dividir un texto en oraciones. La dificultad está en desambiguar el punto.
    \end{block}
    \begin{exampleblock}{Ejemplo}
        ``El Sr. Pérez llegó. ¿Cómo está?'' $\rightarrow$ El punto después de ``Sr.'' no indica fin de oración.
    \end{exampleblock}
    \begin{itemize}
        \item Enfoques: reglas (listas de abreviaturas), modelos ML.
    \end{itemize}
    % Basado en PDF y Pregunta 6
\end{frame}

\begin{frame}{Tokenización por subpalabras}
    \begin{block}{Motivación}
        Resolver el problema de palabras fuera de vocabulario (OOV) y manejar lenguas con morfología rica.
    \end{block}
    \begin{itemize}
        \item \textbf{BPE (Byte-Pair Encoding)}: fusiona iterativamente el par más frecuente.
        \item \textbf{WordPiece}: fusiona el par que maximiza la verosimilitud del lenguaje.
        \item \textbf{Unigram}: modelo probabilístico de subpalabras.
    \end{itemize}
    % Basado en PDF y Pregunta 1
\end{frame}

\begin{frame}{Algoritmo BPE paso a paso}
    \begin{enumerate}
        \item Inicializar vocabulario con caracteres individuales.
        \item Contar frecuencias de pares de símbolos adyacentes.
        \item Fusionar el par más frecuente y añadirlo al vocabulario.
        \item Repetir hasta alcanzar tamaño deseado.
    \end{enumerate}
    \begin{exampleblock}{Ejemplo con ``low low lower''}
        \begin{itemize}
            \item Inicial: l, o, w, e, r
            \item Fusión ``lo'' (frecuente), luego ``low'', etc.
        \end{itemize}
    \end{exampleblock}
    % Basado en PDF y Pregunta 1
\end{frame}

% ============================================================================
% SECCIÓN 3: NORMALIZACIÓN
% ============================================================================
\section{Normalización}

\begin{frame}{Técnicas de normalización}
    \begin{itemize}
        \item \textbf{Minúsculas}: reducir variabilidad (excepto si el caso es informativo).
        \item \textbf{Eliminación de caracteres especiales}: signos de puntuación, caracteres no imprimibles.
        \item \textbf{Manejo de URLs, menciones, hashtags}: reemplazar por tokens especiales (\texttt{<URL>}, \texttt{<USER>}).
        \item \textbf{Manejo de emojis}: eliminar, convertir a texto, o tratar como tokens.
    \end{itemize}
    % Basado en PDF y Pregunta 2
\end{frame}

\begin{frame}[fragile]{Expresiones regulares (regex)}
    \begin{block}{Fundamento matemático}
        Las regex describen lenguajes regulares reconocidos por autómatas finitos deterministas (DFA) en tiempo $O(n)$.
    \end{block}
    \begin{exampleblock}{Ejemplos en Python}
        \begin{lstlisting}
import re
text = re.sub(r'http\S+', '<URL>', text)   # URLs
text = re.sub(r'@\S+', '<USER>', text)     # menciones
text = re.sub(r'#\S+', '<HASHTAG>', text)  # hashtags
text = re.sub(r'[!?.]+', '', text)         # puntuación
        \end{lstlisting}
    \end{exampleblock}
    % Basado en PDF y Pregunta 5
\end{frame}

% ============================================================================
% SECCIÓN 4: ELIMINACIÓN DE STOPWORDS
% ============================================================================
\section{Eliminación de Stopwords}

\begin{frame}{Stopwords: ¿eliminar o no?}
    \begin{itemize}
        \item Palabras muy frecuentes sin mucho significado (artículos, preposiciones).
        \item \textbf{En modelos clásicos (BoW)}: eliminar reduce dimensionalidad y ruido.
        \item \textbf{En modelos contextuales (BERT)}: no es necesario, el modelo aprende a ignorarlas o usarlas para contexto.
        \item \textbf{Cuidado}: palabras como ``no'' son cruciales para sentimiento.
    \end{itemize}
    % Basado en PDF y Pregunta 3
\end{frame}

% ============================================================================
% SECCIÓN 5: STEMMING VS LEMMATIZATION
% ============================================================================
\section{Stemming vs Lemmatization}

\begin{frame}{Stemming}
    \begin{block}{Definición}
        Algoritmo heurístico que elimina sufijos para obtener una raíz (stem). No requiere diccionario.
    \end{block}
    \begin{itemize}
        \item \textbf{Ventaja}: rápido, simple.
        \item \textbf{Desventaja}: produce stems no lingüísticos (``studies'' $\rightarrow$ ``studi'').
        \item Algoritmos: Porter, Snowball, Lancaster.
    \end{itemize}
    % Basado en PDF
\end{frame}

\begin{frame}{Lematización}
    \begin{block}{Definición}
        Reduce a la forma canónica (lema) usando diccionario y contexto (POS tagging).
    \end{block}
    \begin{itemize}
        \item \textbf{Ventaja}: produce palabras reales, más precisa.
        \item \textbf{Desventaja}: más lenta, requiere recursos.
        \item Ejemplo: ``running'' (verbo) $\rightarrow$ ``run''; ``studies'' (sustantivo) $\rightarrow$ ``study''.
    \end{itemize}
    % Basado en PDF
\end{frame}

\begin{frame}{Comparación formal}
    \begin{table}
        \centering
        \begin{tabular}{l|c|c}
            \toprule
            \textbf{Criterio} & \textbf{Stemming} & \textbf{Lematización} \\
            \midrule
            Velocidad & Alta & Baja \\
            Recursos & Ninguno & Diccionario, POS \\
            Precisión lingüística & Baja & Alta \\
            Manejo de irregularidades & No & Sí (``mejor'' $\rightarrow$ ``bueno'') \\
            \bottomrule
        \end{tabular}
    \end{table}
    % Basado en PDF
\end{frame}

\begin{frame}{¿Cuándo usar cada una?}
    \begin{itemize}
        \item \textbf{Stemming}: sistemas de recuperación de información (buscadores) donde la velocidad es crítica.
        \item \textbf{Lematización}: tareas que requieren comprensión semántica (análisis de sentimiento, NER, traducción).
    \end{itemize}
    % Basado en PDF y Pregunta 4
\end{frame}

% ============================================================================
% SECCIÓN 6: OUT-OF-VOCABULARY (OOV)
% ============================================================================
\section{Out-of-Vocabulary (OOV)}

\begin{frame}{El problema OOV}
    \begin{itemize}
        \item Palabras no vistas durante el entrenamiento: neologismos, errores, términos técnicos.
        \item Soluciones tradicionales: token \texttt{<UNK>}, backing-off.
        \item \textbf{Solución moderna: tokenización por subpalabras} (BPE, WordPiece) que descomponen palabras desconocidas en fragmentos conocidos.
    \end{itemize}
    % Basado en PDF y Pregunta 1
\end{frame}

% ============================================================================
% SECCIÓN 7: PIPELINE COMPLETO DE PREPROCESAMIENTO
% ============================================================================
\section{Pipeline Completo de Preprocesamiento}

\begin{frame}{Ejemplo concreto}
    \textbf{Texto original:} ``I LOVE NLP!!! Check out https://example.com for more info @nlp\_expert \#nlp''
    \begin{enumerate}
        \item \textbf{Limpieza con regex}: reemplazar URLs, menciones, hashtags; eliminar puntuación.
        \item \textbf{Minúsculas}.
        \item \textbf{Tokenización} (BPE simulado).
        \item \textbf{Eliminación de stopwords} (opcional).
        \item \textbf{Lematización} (opcional).
    \end{enumerate}
    % Basado en PDF
\end{frame}

\begin{frame}[fragile]{Pipeline paso a paso en código}
    \begin{lstlisting}
import re
text = "I LOVE NLP!!! Check out https://example.com for more info @nlp_expert #nlp"
text = re.sub(r'http\S+', '<URL>', text)
text = re.sub(r'@\S+', '<USER>', text)
text = re.sub(r'#\S+', '<HASHTAG>', text)
text = re.sub(r'[!?.]+', '', text)
text = text.lower()
print(text)
# "i love nlp check out <url> for more info <user> <hashtag>"
    \end{lstlisting}
\end{frame}

% ============================================================================
% SECCIÓN 8: IMPLEMENTACIÓN EN PYTHON CON LIBRERÍAS
% ============================================================================
\section{Implementación en Python}

\begin{frame}[fragile]{NLTK}
    \begin{lstlisting}
import nltk
from nltk.tokenize import word_tokenize
from nltk.corpus import stopwords
from nltk.stem import PorterStemmer, WordNetLemmatizer

nltk.download('punkt')
nltk.download('stopwords')
nltk.download('wordnet')

tokens = word_tokenize("I love NLP. It's fascinating!")
tokens_lower = [t.lower() for t in tokens]
stop_words = set(stopwords.words('english'))
tokens_filtered = [t for t in tokens_lower if t not in stop_words]
stemmer = PorterStemmer()
stems = [stemmer.stem(t) for t in tokens_filtered]
lemmatizer = WordNetLemmatizer()
lemmas = [lemmatizer.lemmatize(t) for t in tokens_filtered]
    \end{lstlisting}
\end{frame}

\begin{frame}[fragile]{spaCy}
    \begin{lstlisting}
import spacy
nlp = spacy.load('en_core_web_sm')
doc = nlp("I love NLP. It's fascinating!")
tokens = [token.text for token in doc]
lemmas = [token.lemma_ for token in doc]
pos_tags = [token.pos_ for token in doc]
    \end{lstlisting}
\end{frame}

\begin{frame}[fragile]{HuggingFace Tokenizers}
    \begin{lstlisting}
from transformers import AutoTokenizer

tokenizer = AutoTokenizer.from_pretrained('bert-base-uncased')
tokens = tokenizer.tokenize("I love NLP!")
input_ids = tokenizer.convert_tokens_to_ids(tokens)
print(tokens, input_ids)
    \end{lstlisting}
\end{frame}

% ============================================================================
% SECCIÓN 9: CASO INTEGRADO - ANÁLISIS DE SENTIMIENTO EN TWITTER
% ============================================================================
\section{Caso Integrado: Análisis de Sentimiento en Twitter}

\begin{frame}{Problema}
    \begin{itemize}
        \item 100,000 tweets etiquetados (positivo/negativo).
        \item Contienen URLs, menciones, emojis, repeticiones de letras (``estoy encantadooo'').
        \item Necesitamos un pipeline de preprocesamiento eficaz.
    \end{itemize}
    % Basado en Pregunta 2
\end{frame}

\begin{frame}{Pasos del pipeline}
    \begin{enumerate}
        \item Reemplazar URLs, menciones, hashtags con tokens especiales.
        \item Normalizar emojis (convertir a texto o token especial).
        \item Reducir repeticiones de letras (ej. ``encantadooo'' $\rightarrow$ ``encantado'').
        \item Minúsculas.
        \item Tokenización por subpalabras (BPE) para manejar jerga.
        \item No eliminar stopwords (``no'' es crucial).
    \end{enumerate}
    % Basado en Pregunta 2
\end{frame}

\begin{frame}{Evaluación}
    \begin{itemize}
        \item Comparar precisión con diferentes pipelines.
        \item Se observa que eliminar URLs mejora precisión; eliminar ``no'' la empeora.
        \item BPE maneja bien palabras como ``bestttt'' (tokenizada como ``best'' + ``ttt'').
    \end{itemize}
\end{frame}

% ============================================================================
% SECCIÓN 10: PREPARACIÓN PARA ENTREVISTAS
% ============================================================================
\section{Preparación para Entrevistas}

\begin{frame}{Preguntas conceptuales típicas}
    \begin{itemize}
        \item ¿Cuándo usar stemming y cuándo lematización?
        \item ¿Qué es BPE y cómo funciona?
        \item ¿Cómo manejas palabras fuera de vocabulario (OOV)?
        \item ¿Siempre eliminas stopwords? ¿Por qué?
    \end{itemize}
    % Basado en PDF
\end{frame}

\begin{frame}{Preguntas prácticas}
    \begin{itemize}
        \item Escribe una expresión regular para eliminar URLs.
        \item ¿Cómo tokenizarías un texto en chino?
        \item ¿Qué impacto tiene la longitud de la secuencia en un Transformer?
    \end{itemize}
    % Basado en Preguntas 1, 5, 12
\end{frame}

% ============================================================================
% SECCIÓN 11: RESUMEN
% ============================================================================
\section{Resumen}

\begin{frame}{Lo que hemos aprendido}
    \begin{exampleblock}{Conceptos clave}
        \begin{itemize}
            \item Tokenización: palabras, oraciones, subpalabras (BPE, WordPiece).
            \item Normalización: minúsculas, regex, manejo de URLs/emojis.
            \item Stopwords: utilidad en modelos clásicos, irrelevante en Transformers.
            \item Stemming vs Lematización: velocidad vs precisión.
            \item OOV resuelto con subword tokenization.
            \item Pipeline completo de preprocesamiento.
            \item Implementaciones con NLTK, spaCy, HuggingFace.
            \item Caso práctico: análisis de sentimiento en Twitter.
        \end{itemize}
    \end{exampleblock}
\end{frame}

% --- Diapositiva final ---
\begin{frame}
    \centering
    \Huge \textbf{¡Gracias!}

\end{frame}

\end{document}